\section{趋势与展望}

下一阶段，国产算力平台上的大模型推理竞争将继续从局部内核优化转向系统能力收敛。图\ref{fig:outlook-map} 将这一判断压缩为四条能力线向同一底座汇聚：真正决定上限的，不再是某个热点算子是否还能再快一些，而是平台边界、运行时组织、统一评测与交付治理能否同时稳定。换言之，未来的差距主要体现在谁能把异构硬件、复杂负载和快速演进的模型生态压回同一套可持续的工程框架，而不是把问题分散为若干彼此割裂的性能补丁。

\begin{figure}[t]
\centering
\resizebox{0.98\linewidth}{!}{\begin{tikzpicture}[
    node distance=18mm and 24mm,
    block/.style={draw=black!75, line width=0.8pt, rounded corners=2.2pt, minimum width=31mm, minimum height=10mm, align=center, font=\small, inner sep=4pt},
    centerblock/.style={draw=black!75, line width=0.8pt, rounded corners=3pt, minimum width=40mm, minimum height=12mm, align=center, font=\small\bfseries, inner sep=5pt},
    note/.style={font=\footnotesize, align=center, fill=white, inner sep=1.2pt, text=black!80},
    flow/.style={-{Latex[length=2.2mm]}, semithick, draw=black!75}
]
\node[centerblock, fill=gray!8] (goal) {可持续的国产算力\\推理系统底座};

\node[block, fill=blue!9, above left=20mm and 28mm of goal] (plugin) {后端插件化\\平台抽象/注册/可合并优化};
\node[block, fill=green!10, above right=20mm and 28mm of goal] (runtime) {资源管理与调度\\SLO/迁移/KV Cache/阶段解耦};
\node[block, fill=orange!12, below left=20mm and 28mm of goal] (eval) {统一评测与观测\\近真实负载/成本/稳定性};
\node[block, fill=red!8, below right=18mm and 24mm of goal] (model) {模型结构持续演进\\MoE/MLA/多模态/语音交互};

\draw[flow] (plugin.south east) -- (goal.north west);
\draw[flow] (runtime.south west) -- (goal.north east);
\draw[flow] (eval.north east) -- (goal.south west);
\draw[flow] (model.north west) -- (goal.south east);

\node[note, text width=6.2cm, above=6mm of goal] {未来竞争将从单点热点优化转向系统能力收敛。};
\node[note, text width=6.4cm, below=6mm of goal] {平台抽象、运行时、评测与模型演进需要同步推进。};
\end{tikzpicture}}
\caption{国产推理引擎未来演进方向关系图}
\label{fig:outlook-map}
\end{figure}

首先会成为分水岭的，是平台抽象能否从“后端接入机制”演进为“能力契约”。随着上游框架持续扩展多模态、结构化输出和 Agent 场景，国产后端若仍依赖侵入式 patch 和深度分叉，就很难长期吸收社区演进成果。更有持续价值的做法，是把平台差异收敛为可合并、可替换、可验证的模块，并把 attention backend、图模式、动态 shape、通信语义、低比特路径和前处理约束显式化，使调度器、模型注册器和服务层能够基于能力而不是经验特判做决策。对本土项目而言，平台插件化若能进一步沉淀为稳定契约，其价值会明显高于单轮适配带来的局部性能收益。

其次，围绕长上下文、多模态与复杂控制流的资源管理会成为运行时优化中心。Prefix Cache、分层 KV Cache、跨实例迁移、SLO 感知调度和多模态阶段解耦，本质上都在处理同一个问题：如何在不确定负载下稳定组织状态，而不是被状态牵着走。Llumnix、SOLA、LServe、AI Metropolis 与 ElasticMM 等工作已经说明，后续优化重点会更多落在状态迁移、依赖感知执行、统一稀疏注意力和多模态弹性并行等机制上\citep{sun2024llumnix,hong2025sola,yang2025lserve,xie2025aimetropolis,liu2025elasticmm}。对国产平台而言，这一点尤为关键，因为图编译成本、shape 敏感性和设备差异往往会把局部调度问题直接放大为端到端服务问题。未来更稳健的实现方向，不是继续把新能力外挂到 token 循环之外，而是把请求状态、缓存状态、约束状态与回退状态纳入统一状态机。

第三，压缩策略与成本核算会从离线配置问题转向在线协同问题。未来若要真正降低单位 token 成本，系统不能只记录某组量化参数在离线评测中的精度损失，还必须同时知道它会怎样改变图模式命中、状态驻留、迁移带宽与尾时延。谁能更早把混合精度、KV 压缩、稀疏化路径和运行时观测组织成统一收益--代价闭环，谁才更有可能把“低成本”沉淀为稳定能力，而不是停留在阶段性实验结果上。

最后，统一评测与交付控制面将决定这些优化能否沉淀为真实生产力。不同后端只有在统一接口和近真实负载下被比较，优化方向才不会被展示型 benchmark 带偏；部署环境中的驱动、编译链、插件版本和依赖树只有被纳入自动化诊断、回归和版本矩阵校验，系统上线与回滚成本才不会持续上升。vLLM-HUST 将平台插件、运行时护栏和外部环境治理拆分组织的做法，已经体现出“引擎主体 + 治理配套工具链”的演进方向\citep{vllmhust}。这也与 2025--2026 年公开社区高频出现的版本对齐、PD 分离、CANN/torch\_npu 兼容和 RoCE/HCCL 连通性问题相一致\citep{zhihu_vllm_ascend_guide,zhihu_vllm_pd_migration}。

模型侧的变化还会继续重写推理系统的受力点。DeepSeek-V2/V3 所代表的 MLA 与 MoE 路线，正在把通信组织、缓存布局和专家路由重新推回运行时中心位置\citep{liu2024deepseekv2,deepseekv3}；Qwen2-VL、InternVL2、CogVLM2、GLM-4-Voice 与 MiniCPM-o 等模型则说明，视觉、视频和语音链路已经使文本时代的单一路径假设难以维持\citep{wang2024qwen2vl,chen2024internvl2,wang2024cogvlm2,zeng2024glm4voice,cui2026minicpmo45}。因此，下一阶段最值得警惕的不是“还有哪些优化尚未完成”，而是系统是否仍沿用面向单模型、单平台、单 benchmark 的局部竞赛逻辑。只有把平台契约、运行时状态机、评测闭环和交付治理一起推进，国产推理引擎才能从阶段性适配进入可持续积累。
